\section{Timing And Further Reading}\label{sec:further-reading}

Fourier prepares each spectrum over one Hop interval and displays it when
complete. Longer Length settings include more past audio, Average slows the
response further, and Rack's display refresh adds its own delay. A shorter
Hop therefore gives more frequent analyses without making a long frame
instantaneous.

Window, smoothing, frequency bounds, scales, and Slope are applied to the
next analysis frame. Changing Length clears captured audio; changing Rack's sample rate
also discards unfinished analysis. Let new audio fill the frame before making
comparisons after either change. Fourier continues analyzing retained audio
when Run is off, as described in the Run section.

Analysis runs on Rack's audio engine thread. Work is spread across sample
calls to reduce processing bursts, but the module still contributes to the
patch's CPU load. This does not guarantee that a heavily loaded patch will
meet its audio deadlines.

For FFT scheduling, the analysis method, and performance experiments, see
the accompanying
\href{https://github.com/Kautenja/ArhythmeticUnits-Fourier/tree/main/docs/whitepaper}{\textit{Scheduling FFT-Based Spectral Analysis in the Audio Processing Loop}}.
Detailed proofs and supporting algorithms are retained separately with the
paper; this manual focuses on operating the module.
